\subsection{锥；子空间的缩塌}

设 X 与 Y 是非空拓扑空间，\(f: X \to Y\) 是连续映射。令 \(\operatorname{Côn}(f)\) 为映射 \(f\) 的锥，\(s\) 为其顶点。用 \(\alpha_{f}^{\prime}: X \times I \to \operatorname{Côn}(f)\) 与 \(\beta_{f}^{\prime}: Y \to \operatorname{Côn}(f)\) 表示规范映射。\(\alpha_{f}^{\prime}\) 在子空间 \(X \times \{0\}\)（\(X \times I\) 的）上的限制是以 \(\{s\}\) 为像的常值映射。映射 \(\beta_{f}^{\prime}\) 诱导出 Y 到锥 \(\operatorname{Côn}(f)\) 的锥底上的一个同胚，我们将借此把这两个空间等同起来。再记

\[
\sigma_ {f} ^ {\prime} \colon (\operatorname{Côn} (f) - \{s \}) \times \mathbf {I} \to \operatorname{Côn} (f) - \{s \}
\]

为规范收缩，并记 \(\rho_{f}^{\prime}\colon\operatorname{Côn}(f) - \{s\}\to\mathrm{Y}\) 为去掉顶点的锥到其锥底上的规范收缩映射。

令 \(J = \pi_{0}(X)\)；对 J 的每个元素 j，用 \(X_{j}\) 表示 X 的第 j 个分支，即令 \(b_{j}\) 为 \(X_{j}\) 的一点，并用 \(\gamma_{j}\) 表示道路 \(t \mapsto \alpha_{f}^{\prime}(b_{j}, t)\)（在 \(\operatorname{Côn}(f)\) 上，连接 \(s\) 与 \(f(b_{j})\)）的类。

令 I 为映射 \(\pi_{0}(f)\) 的像：这是 Y 的与 \(f(\mathrm{X})\) 相遇的道路连通分支组成的集合；用 \(\varphi\colon J\to I\) 表示由 f 过渡到道路连通分支而导出的映射。对每个元素 \(i\in I\)，用 \(Y_{i}\) 表示 Y 的分支 i，并选取一点 \(a_{i}\)（在 \(Y_{i}\) 上）。

对 J 的每个元素 \(j\)，选取一条道路 \(\mathrm{B}_{j}\)，它连接点 \(f(b_{j})\) 与点 \(a_{\varphi(j)}\)（在 \(\mathrm{Y}_{\varphi(j)}\) 上），并用 \(\beta_{j}\) 表示它的类。用 \(\psi_{j}\) 表示从 \(\pi_{1}(\mathrm{X},b_{j})\) 到 \(\pi_{1}(\mathrm{Y},a_{\varphi(j)})\) 中由

\[
\psi_ {j} (v) = \beta_ {j} ^ {- 1} f _ {*} (v) \beta_ {j}
\]

所定义的同态，其中 \(v \in \pi_1(X, a_j)\)。

设 \(\sigma\colon I\to J\) 是映射 \(\varphi\) 的一个截面。令 \(T=\sigma(I)\)，\(\tau=\sigma\circ\varphi\)；注意 \(\varphi\circ\tau=\varphi\)。

命题 4. --- 假设 Y 的道路连通分支都是开集。对每个 \(i \in I\)，令 \(G_i\) 为群 \(\pi_1(Y_i, b_i)\) 关于包含诸同态 \(\psi_j\)（\(j \in \overline{\varphi}^{-1}(i)\)）的像的最小正规子群所得的商；并用 \(p_i\) 表示从 \(\pi_1(Y_i, b_i)\) 到 \(G_i\) 上的规范满射。

存在唯一的群同态

\[
\mathrm{P} \colon \underset {i \in \mathrm{I}} {*} \mathrm{G} _ {i} * \mathrm{F} (\mathrm{J} - \mathrm{T}) \to \pi_ {1} (\operatorname{C} \hat {\mathrm{n}} (f), s)
\]

使得

\[
\begin{array}{c c} \mathsf {P} (p _ {i} (v)) = \operatorname{Int} (\gamma_ {\sigma (i)} \beta_ {\sigma (i)}) (v) & \text {for } i\in I \text { and } v\in\pi_{1} (Y_{i},b_{i}), \\ \mathsf {P} (j) = \gamma_ {j} \beta_ {j} \beta_ {\tau (j)} ^ {- 1} \gamma_ {\tau (j)} ^ {- 1} & \text {for } j\in J - T. \end{array}
\]

同态 P 是同构。

用 Y' 表示 Y 的与 \(f(\mathrm{X})\) 相遇的道路连通分支的并，并令 \(f'\colon X \to Y'\) 为由 \(x \mapsto f(x)\) 给出的映射。集合 Y' 是 Y 的开子集；其道路连通分支都是开集。锥 \(\operatorname{Côn}(f')\) 等同于 s 在 \(\operatorname{Côn}(f)\) 中的道路连通分支。这使我们可以假设 \(Y = Y'\)，换言之，映射 \(\pi_{0}(f)\) 是满射的且 \(I = \pi_{0}(Y)\)。

对每个 \(j \in \mathrm{J}\)，令 \(\mathrm{V}_j = \alpha_f'(X_j \times ]0,1[)\)。通过过渡到子空间，映射 \(\alpha_f'\) 诱导出 \(X \times ]0,1[\) 到 \(Y \cup \{s\}\) 在 \(\operatorname{Côn}(f)\) 中的补集上的同胚。于是，诸集合 \(V_j\) 是 \(\operatorname{Côn}(f) - (Y \cup \{s\})\) 的道路连通分支。

对每个 \(i \in \mathrm{I}\)，令 \(\mathrm{U}_i = (\rho_f')^{-1}(\mathrm{Y}_i)\)；它是 \(\operatorname{Côn}(f)\) 的开子集，因为由假设 \(\mathrm{Y}_i\) 在 \(\mathrm{Y}\) 中开。对每个 \(j \in \overline{\varphi}^{-1}(i)\)，我们有 \(f(\mathrm{X}_j) \subset \mathrm{Y}_i\)，且

\[
\mathrm{V} _ {j} \cup \mathrm{Y} _ {i} = \alpha_ {f} ^ {\prime} (\mathrm{X} _ {j} \times ] 0, 1 ]) \cup \mathrm{Y} _ {i},
\]

于是 \(V_{j} \cup Y_{i}\) 是 \(\operatorname{Côn}(f)\) 的包含 \(Y_{i}\) 的道路连通子集。由于 \(U_{i}\) 是 \(Y_{i}\) 与诸集合 \(V_{j}\)（\(j \in \overline{\varphi}^{-1}(i)\)）的并，由此推出 \(U_{i}\) 是道路连通的。

最后，集合 \(\mathrm{C}'(\mathrm{X}) = \mathrm{C}\hat{\mathrm{o}}\mathrm{n}(f) - \mathrm{Y}\) 是 \(\mathrm{C}\hat{\mathrm{o}}\mathrm{n}(f)\) 的开子集；它可收缩到 s，因此是道路连通的。

集合 C'(X) 与诸集合 \(U_{i}\)（\(i \in I\)）构成 \(\operatorname{Côn}(f)\) 的一个由开道路连通子集组成的覆盖，我们将对它应用 IV, 第~412~页的命题 1。令 I' 为在 I 中添上 s 所得的集合；给它赋以使 s 为其最小元素的全序。

对 I 的不同元素 \(i, i'\)，有 \(\mathrm{U}_{i} \cap \mathrm{U}_{i'} = \varnothing\)。对 \(i \in \mathrm{I}\)，\(\mathrm{C}'(\mathrm{X}) \cap \mathrm{U}_{i}\) 是诸集合 \(\mathrm{V}_{j}\)（\(j \in \overline{\varphi}^{-1}(i)\)）的并；它们道路连通且两两不交。该覆盖的任意三个不同集合的交都是空集。

所考虑覆盖的框架 \(\Gamma\) 以集合 I' 为顶点集。其箭为诸三元组 \((s, i, V_j)\)，其中 \(j \in \mathbf{J}\)，\(i = \varphi(j)\)；这样我们就把 \(\Gamma\) 的箭集等同于集合 J。

对 \(i \in I\)，我们取基点 \(\mathsf{a}(i) = a_i \in \mathrm{U}_i\)；并令 \(\mathsf{a}(s) = s \in \mathrm{C}'(\mathrm{X})\)。

对 \(j \in J\)，令 \(\mathsf{b}(j) = \alpha_f'(b_j, \frac{1}{2})\)。用 \(B_1(j)\) 表示 \(C'(X)\) 中以 \(\mathsf{b}(j)\) 为源、以 \(\mathsf{a}(s)\) 为终点、由 \(t \mapsto \alpha_{f}^{\prime}(b_{j}, (1 - t)/2)\) 给出的道路。用 \(\mathrm{B}_{2}(j)\) 表示 \(\mathrm{U}_{\varphi(j)}\) 中以 \(\mathfrak{b}(j)\) 为源、以 \(\mathfrak{a}(\varphi(j)) = a_{\varphi(j)}\) 为终点、由道路 \(t \mapsto \alpha_{f}^{\prime}(b_{j}, (1 + t)/2)\) 与道路 \(\mathrm{B}_{j}\) 的拼接给出的道路。于是，道路 \(\mathrm{B}(j) = \overline{\mathrm{B}}_{1}(j) * \mathrm{B}_{2}(j)\) 的类等于 \(\gamma_{j}\beta_{j}\)。

我们取 \(i_{0}=s\)。

我们取 \(\Gamma\) 的唯一以 \(\sigma(\mathrm{I})\) 为箭集的定向树作为极大定向树 T。我们有 \(\delta(s) = e\)，而对 \(i \in \mathrm{I}\)，\(\delta(i) = [\mathrm{B}(\sigma(i))] = \gamma_{\sigma(i)}\beta_{\sigma(i)}\)。

设 \(j \in \mathrm{J}\)，且 \(i = \varphi(j)\)。

同态 \(\varphi_{j}\)（从 \(\pi_1(V_j, \mathfrak{b}(j))\) 到 \(\pi_1(C'(X), s)\) 中）是平凡同态，因为 \(C'(X)\) 可收缩到 s。

映射 \(\alpha_{f}^{\prime}\) 诱导出 \(\mathrm{X}_{j} \times ]0,1[\) 到 \(\mathrm{V}_{j}\) 上的一个同胚；这个同胚诱导出群 \(\pi_{1}(\mathrm{V}_{j},\mathsf{b}(j))\) 到群 \(\pi_{1}(\mathrm{X}_{j},b_{j}) = \pi_{1}(\mathrm{X},b_{j})\) 上的同构。通过过渡到子空间，映射 \(\sigma_{f}^{\prime}\) 诱导出 \(\mathrm{U}_{i}\) 到 \(\mathrm{Y}_{i}\) 上的一个强收缩，它诱导出群 \(\pi_{1}(\mathrm{U}_{i},\mathsf{a}(i))\) 到群 \(\pi_{1}(\mathrm{Y},a_{i})\) 上的同构。借助于这些同构，同态

\[
\psi_ {j} \colon \pi_ {1} (\mathrm{V} _ {j}, \mathsf {b} (j)) \to \pi_ {1} (\mathrm{U} _ {\varphi (j)}, \mathsf {a} (\varphi (j)))
\]

就等同于同态 \(\operatorname{Int}(\beta_{\varphi(j)}^{-1}) \circ f_*\)（从 \(\pi_1(\mathrm{X}, b_j)\) 到 \(\pi_1(\mathrm{Y}, a_i)\) 中）。

由于 C'(X) 可收缩到 s，同态 \(\mu_{s}\) 是平凡同态。

设 \(i \in I\)。同态 \(\mu_i\)（从 \(\pi_1(U_i, \mathfrak{a}(i))\) 到 \(\pi_1(\operatorname{Côn}(f), s)\) 中）被等同于这样的同态：它由从 \(\pi_1(Y, a_i)\) 到 \(\pi_1(\operatorname{Côn}(f), s)\) 中的同态 \(\operatorname{Int}(\delta(i))\) 与从 \(\pi_1(U_i, s)\) 到 \(\pi_1(\operatorname{Côn}(f), s)\) 中的同态（由 \(U_i\) 到 \(\operatorname{Côn}(f) \) 的规范单射导出）复合而成。

最后，同态 \(\mu\colon\mathrm{F}(\mathrm{J})\to\pi_{1}(\mathrm{C}\hat{\mathrm{o}}\mathrm{n}(f),s)\) 由下式给出

\[
\mu (j) = \gamma_ {j} \beta_ {j} \beta_ {\tau (j)} ^ {- 1} \gamma_ {\tau (j)} ^ {- 1}.
\]

令 P' 为唯一的群同态

\[
\mathrm{P} ^ {\prime} \colon \underset {i \in \mathrm{I}} {*} \pi_ {1} (\mathrm{Y} _ {i}, a _ {i}) * \mathrm{F} (\mathrm{J}) \to \pi_ {1} (\mathrm{C} \hat {\mathrm{o}} \mathrm{n} (f), s)
\]

它与同态 \(\operatorname{Int}(\delta(i)^{-1})\) 在 \(\pi_1(Y_i, a_i)\) 上、与同态 \(\mu\) 在 \(\mathrm{F}(\mathrm{J})\) 上一致。由 IV, 第~412~页的命题 1，同态 P' 是满射的，其核是 \(*\pi_1(Y_i, a_i) * F(J)\) 中包含诸元素 \(r_1(j)\)（\(j \in T\)）以及诸元素 \(r_{2}(j, v)\)（\(j \in J\) 且 \(v \in \pi_{1}(V_{j}, b(j))\)）的最小正规子群。（由于集合 K 是空集，不存在元素 \(r_{3}(k)\)。）

对 \(j \in \mathrm{T}\)，我们有 \(r_1(j) = j\)。设 \(j \in \mathrm{J}\)，并设 \(v \in \pi_1(\mathrm{V}_j, \mathfrak{b}(j))\)。把 \(\pi_1(\mathrm{V}_j, \mathfrak{b}(j))\) 与 \(\pi_1(\mathrm{X}, b_j)\) 的等同考虑在内，我们有 \(r_2(j, v) = j\psi_j(v)j^{-1}\)。用 \(p\) 表示从 \(*_{i}\pi_1(\mathrm{Y}_i, a_i) * \mathrm{F}(\mathrm{J})\) 到 \(*_{i}\mathrm{G}_{i} * \mathrm{F}(\mathrm{J} - \mathrm{T})\) 上的规范满射同态。对 \(j \in \mathrm{T}\)，我们有 \(p(r_1(j)) = e\)；对 \(j \in \mathrm{J}\) 与 \(v \in \pi_1(\mathrm{X}, b_j)\)，我们有 \(p(\psi_j(v)) = e\)。因此，存在唯一的群同态 \(\mathsf{P}\)，从 \(*_{i}\mathrm{G}_{i} * \mathrm{F}(\mathrm{J} - \mathrm{T})\) 到 \(\pi_1(\operatorname{Côn}(f), s)\) 中，使得 \(\mathsf{P}' \circ p = \mathsf{P}\)；它是同构。

推论 1. --- 再设空间 X 与 Y 都道路连通，并设 a 是 X 的一点。从 \(\pi_1(Y, f(a))\) 到 \(\pi_1(\operatorname{Côn}(f), f(a))\) 的规范映射是满射的，其核是包含同态 \(\pi_1(f, a)\) 的像的最小正规子群。

推论 2. --- 再设 Y 的道路连通分支都是道路单连通的。那么同态 \(\mu\colon F(J-T)\to\pi_{1}(\operatorname{Côn}(f),s)\) 是同构。

注 1. --- 设 X 是道路连通分支均为开集的拓扑空间。设 A 是 X 的闭子空间；用 \(\iota\colon A\to X\) 表示规范单射，用 X/A 表示把 A 缩塌为一点后由 X 得到的空间，并用 \(o\) 表示 A 被缩塌到的那一点，用 \(p\colon X\to X / A\) 表示规范映射。再设偶 (X,A) 具有同伦扩张性质。于是规范映射 \(\overline{\rho}\colon \operatorname{Côn}(\iota)\to \mathrm{X} / \mathrm{A}\) 是同伦等价（III, 第~255~页，注 1），而由命题 4 可得出 X/A 在其基点 \(o\) 处的庞加莱群的计算。特别地，若 X 的道路连通分支都是道路单连通的，则群 \(\pi_1(\mathrm{X / A},o)\) 是自由群。
% semantic-fragments: legacy-input-seam
\relax{}
